\section{Repository Integration and Chapter Dependencies}
\label{sec:ch4-integration}

The vocabulary of this chapter is now part of the book-wide glossary and notation system. In particular, \gls{ordinarydifferentialequation}, \gls{initialvalueproblem}, \gls{directionfield}, \gls{phasespace}, \gls{equilibrium}, \gls{stability}, \gls{bifurcation}, and \gls{stiffness} are used consistently in later chapters. The canonical state-space convention is
\[
\dot{\mathbf x}=\mathbf f(t,\mathbf x),
\qquad
\mathbf x(t_0)=\mathbf x_0,
\]
while autonomous systems omit the explicit time argument.

\begin{physicalbox}[Dependency map: from laws of change to physics]
\[
\begin{array}{c}
\text{initial-value problem}\longrightarrow\text{trajectory}\longrightarrow\text{phase flow}\\[2mm]
\downarrow\hspace{30mm}\downarrow\hspace{30mm}\downarrow\\[-1mm]
\text{Newtonian motion}\qquad\text{circuit dynamics}\qquad\text{quantum evolution}
\end{array}
\]
Chapter~\ref{ch:mathematical-language-of-physics} supplies vectors and eigenmodes; Chapter~\ref{ch:differential-calculus} supplies linearization and Jacobians; Chapter~\ref{ch:vector-calculus} supplies spatial field operators. This chapter combines them into evolution laws. Later mechanics chapters use second-order equations and phase space, electromagnetism uses coupled first-order field evolution, and quantum mechanics uses linear state evolution generated by a Hamiltonian.
\end{physicalbox}

\index{ordinary differential equation}
\index{initial-value problem}
\index{direction field}
\index{phase space}
\index{equilibrium point}
\index{stability}
\index{bifurcation}
\index{stiffness}
\index{Runge--Kutta method}
\index{harmonic oscillator}
